\notechapter{N-20230123}{Definite Integrals, Line Integrals, Work, Green-Type Thinking, and Vector Operations}

\section{Indefinite and definite integrals}

Indefinite and definite integrals play different roles.  An indefinite integral is an antiderivative family:
\[
  \int f(x)\,dx=F(x)+C,\qquad F'(x)=f(x).
\]
A definite integral is a number:
\[
  \int_a^b f(x)\,dx.
\]
Geometrically, it is signed area under the graph.  Analytically, it is the limit of Riemann sums.

The fundamental theorem of calculus connects the two:
\[
  \int_a^b f(x)\,dx=F(b)-F(a)
\]
when $F'=f$.

\section{Area by slicing}

The page includes diagrams of shaded regions.  The standard method is to slice a region into thin strips.  If vertical strips are convenient, then
\[
  A=\int_a^b (y_{top}(x)-y_{bottom}(x))\,dx.
\]
If horizontal strips are more convenient, then
\[
  A=\int_c^d (x_{right}(y)-x_{left}(y))\,dy.
\]
Choosing the direction of slicing is part of the solution, not a fixed rule.

\section{Work as an integral}

The physical meaning of a definite integral appears in work.  If a force $F(x)$ acts along a line, then the work from $x=a$ to $x=b$ is
\[
  W=\int_a^b F(x)\,dx.
\]
If the force is constant, this reduces to $W=F(b-a)$.  If the force changes with position, integration adds the small contributions
\[
  dW=F(x)\,dx.
\]
This is the same accumulation idea as area, but with a physical interpretation.

\section{Line integrals}

The note then introduces line integrals.  If a curve is parametrized by
\[
  r(t)=(x(t),y(t)),\qquad a\le t\le b,
\]
then the scalar line element is
\[
  ds=\sqrt{(x'(t))^2+(y'(t))^2}\,dt.
\]
A scalar line integral is
\[
  \int_C f\,ds
  =\int_a^b f(x(t),y(t))\sqrt{(x'(t))^2+(y'(t))^2}\,dt.
\]
For a vector field $F=(P,Q)$, the work integral is
\[
  \int_C F\cdot dr
  =\int_C P\,dx+Q\,dy
  =\int_a^b [P(x(t),y(t))x'(t)+Q(x(t),y(t))y'(t)]\,dt.
\]

\section{Green-type intuition}

An oriented boundary and a vector field suggest that circulation around the boundary should be related to a double integral over the region.  In the plane, Green's theorem says
\[
  \oint_{\partial D} P\,dx+Q\,dy
  =
  \iint_D \left(\frac{\partial Q}{\partial x}-\frac{\partial P}{\partial y}\right)\,dA.
\]
Even if the note does not prove the theorem fully, it records the key intuition: local rotation accumulated over an area produces circulation on the boundary.

\section{Polar coordinates}

In polar coordinates, the change of variables is
\[
  x=r\cos\theta,\qquad y=r\sin\theta.
\]
The area element becomes
\[
  dA=r\,dr\,d\theta.
\]
The factor $r$ is the Jacobian.  Geometrically, a small polar rectangle has side lengths approximately $dr$ and $r\,d\theta$.

Thus a disk of radius $R$ has area
\[
  \int_0^{2\pi}\int_0^R r\,dr\,d\theta=\pi R^2.
\]

\section{Spherical coordinates}

For spherical coordinates, a common convention is
\[
  x=\rho\sin\varphi\cos\theta,\qquad
  y=\rho\sin\varphi\sin\theta,\qquad
  z=\rho\cos\varphi.
\]
The volume element is
\[
  dV=\rho^2\sin\varphi\,d\rho\,d\varphi\,d\theta.
\]
The factor $\rho^2\sin\varphi$ comes from the three small lengths:
\[
  d\rho,\qquad \rho\,d\varphi,\qquad \rho\sin\varphi\,d\theta.
\]
This explains the diagram on the final page: the small spherical patch has dimensions determined by radius and angle.

\section{Dot product and cross product}

The note also summarizes vector operations.  The dot product is
\[
  a\cdot b=|a||b|\cos\theta.
\]
It measures projection and appears in work:
\[
  dW=F\cdot dr.
\]
The cross product has magnitude
\[
  |a\times b|=|a||b|\sin\theta
\]
and direction given by orientation.  It measures the area of the parallelogram spanned by $a$ and $b$.

In coordinates,
\[
  a\times b=
  \begin{vmatrix}
  i&j&k\\
  a_1&a_2&a_3\\
  b_1&b_2&b_3
  \end{vmatrix}.
\]

\section{After the examples: what object is being integrated?}

The examples develop integration as accumulation over different geometric objects.  The clean way to organize them is not by dimension alone, but by the type of object being integrated.

\begin{center}
\small
\begin{tabularx}{\textwidth}{@{}lYY@{}}
\toprule
Integral & Object being integrated & Geometric meaning\\
\midrule
\(\int_a^b f(x)\,dx\) & density on an oriented interval & signed accumulation\\
\(\int_C f\,ds\) & scalar density along a curve & mass/length-weighted accumulation\\
\(\int_C P\,dx+Q\,dy\) & one-form & work or circulation\\
\(\iint_D f\,dA\) & area density & accumulation over a region\\
\(\iiint_\Omega f\,dV\) & volume density & weighted volume or mass\\
\bottomrule
\end{tabularx}
\end{center}

This table prevents a common confusion.  A line integral of \(f\,ds\) and a work integral \(P\,dx+Q\,dy\) both run along a curve, but they integrate different geometric objects.  The first uses arclength and ignores orientation; the second is a pullback of a one-form and changes sign when the orientation is reversed.

Green's theorem then becomes a boundary-to-interior conversion:
\[
\oint_{\partial D}P\,dx+Q\,dy
=
\iint_D(\partial_x Q-\partial_y P)\,dA.
\]
It says that circulation around the boundary is curl accumulated over the interior.

\section{Line integrals as pullbacks}

For a parametrized curve $\gamma:[a,b]\to\mathbb R^n$, a line integral of a one-form
\[
  \omega=P_1dx_1+\cdots+P_ndx_n
\]
is
\[
  \int_\gamma \omega
  =\int_a^b \sum_i P_i(\gamma(t))\gamma_i'(t)\,dt.
\]
This is the pullback of the one-form to the interval:
\[
  \gamma^*\omega=\sum_i P_i(\gamma(t))\gamma_i'(t)dt.
\]
Thus the elementary formula $Pdx+Qdy$ is already differential-form notation.

\section{Exact forms and potentials}

A one-form $\omega$ is exact if $\omega=df$ for some potential $f$.  Then
\[
  \int_\gamma df=f(\gamma(b))-f(\gamma(a)).
\]
This is the fundamental theorem of line integrals.

A form is closed if $d\omega=0$.  In the plane, for $\omega=Pdx+Qdy$, this means
\[
  Q_x-P_y=0.
\]
On simply connected domains, closed one-forms are exact.  On domains with holes, this may fail.

\section{Conservation laws}

In mechanics, work is a line integral
\[
  W=\int_C F\cdot dr.
\]
If $F=-\nabla V$ is conservative, then
\[
  W=V(A)-V(B)
\]
depends only on endpoints.  Conservation of energy is therefore a statement about exact one-forms.

This connects elementary work integrals to topology: path independence is not only about derivatives; it also depends on whether the domain has holes.

\section{De Rham preview}

De Rham cohomology measures the failure of closed forms to be exact:
\[
  H^1_{dR}(M)=\frac{\{\text{closed 1-forms}\}}{\{\text{exact 1-forms}\}}.
\]
For the punctured plane, the form
\[
  \omega=\frac{-y\,dx+x\,dy}{x^2+y^2}
\]
is closed but not exact.  Its integral around the unit circle is $2\pi$.

The elementary line integral around a circle is therefore a measurement of a topological hole.

\section{Work integrals and gauge freedom}

If a force field is conservative, $F=-\nabla V$, then work depends only on endpoints.  The potential $V$ is not unique: adding a constant changes nothing.  This is the simplest form of gauge freedom.  The physically meaningful object is the differential $dV$, not the absolute zero of potential.

In electromagnetism, potentials have richer gauge freedom.  A vector potential $A$ may be changed by a gradient without changing the magnetic field $B=\nabla\times A$.  The elementary fact that potentials are determined only up to constants is the first example of this broader principle.

\section{Polar and spherical adapted charts}

The original integral examples using polar or spherical coordinates should be read as choosing charts adapted to symmetry.  A chart is a coordinate map from a parameter domain to the geometric region.  The integration element is obtained by pulling back the metric or volume form.

Thus polar coordinates are not merely a trick for circular regions.  They are a coordinate chart on the punctured plane adapted to the action of rotations.  Spherical coordinates are adapted to radial symmetry in three dimensions.
